\section*{अध्याय \ref{ch:g12:ka-and-pka} --- अम्ल और क्षारक: Ka और pKa}

\begin{solution}{exo:g12:ka-and-pka:1}
संयुग्मी क्षारक: \ce{HCOO-}, \ce{NH3}, \ce{OH-}, \ce{CO3^{2-}}। संयुग्मी
अम्ल: \ce{NH4+}, \ce{H2O}, \ce{H3O+}, \ce{CO2,H2O}। \omterm{def:g12:ka-and-pka:oxonium}{उभयधर्मी}: \ce{H2O}
और \ce{HCO3-}।
\end{solution}

\begin{solution}{exo:g12:ka-and-pka:2}
\omterm{def:g9:acids-bases-ph:ph}{pH} 1.0; 3.0; $-\log(\num{5.0e-3}) = 2.30$।
\end{solution}

\begin{solution}{exo:g12:ka-and-pka:3}
$10^{-4.5} = \qty{3.2e-5}{mol/L}$। \omterm{def:g9:acids-bases-ph:ph}{pH} 11.2 पर:
$[\ce{H3O+}] = \qty{6.3e-12}{mol/L}$ और
$[\ce{OH-}] = \num{1.0e-14}/\num{6.3e-12} = \qty{1.6e-3}{mol/L}$।
\end{solution}

\begin{solution}{exo:g12:ka-and-pka:4}
\ce{HCOOH + H2O <=> HCOO- + H3O+},
$K_a = \dfrac{[\ce{HCOO-}][\ce{H3O+}]}{[\ce{HCOOH}]\,c^\circ}$;
\ce{NH4+ + H2O <=> NH3 + H3O+},
$K_a = \dfrac{[\ce{NH3}][\ce{H3O+}]}{[\ce{NH4+}]\,c^\circ}$।
\end{solution}

\begin{solution}{exo:g12:ka-and-pka:5}
मेथेनोइक अम्ल (छोटा $pK_a$)। अधिक \omterm{def:g12:ka-and-pka:strong-acid}{प्रबल क्षारक} \ce{CH3COO-} है, जो
अधिक \omterm{def:g12:ka-and-pka:strong-acid}{दुर्बल अम्ल} का क्षारक है।
\end{solution}

\begin{solution}{exo:g12:ka-and-pka:6}
$[\ce{A-}] = [\ce{H3O+}] = \qty{1.0e-3}{mol/L}$;
$[\ce{HA}] = 0.050 - 0.001 = \qty{0.049}{mol/L}$;
$K_a = \num{1.0e-6}/0.049 = \num{2.0e-5}$; $pK_a = 4.69$।
\end{solution}

\begin{solution}{exo:g12:ka-and-pka:7}
$K_a = \num{6.31e-5}$; $x^2 + \num{6.31e-5}\,x - \num{1.26e-6} = 0$ से
$x = \qty{1.09e-3}{mol/L}$ मिलता है: \omterm{def:g9:acids-bases-ph:ph}{pH} 2.96।
\end{solution}

\begin{solution}{exo:g12:ka-and-pka:8}
$\mathrm{pH} = 14.0 + \log(\num{2.0e-3}) = 14.0 - 2.70 = 11.3$।
\end{solution}

\begin{solution}{exo:g12:ka-and-pka:9}
\ce{HF}, \ce{HCOOH}, ऐस्कॉर्बिक अम्ल, \ce{C6H5COOH}, \ce{CH3COOH},
\ce{CO2,H2O}, \ce{NH4+}, \ce{HCO3-}। \omterm{def:g12:ka-and-pka:strong-acid}{प्रबल अम्ल} बहुत नीचे,
पैमाने के छोर से भी परे, होता है: वह पानी से पूर्ण रूप से अभिक्रिया करता है।
\end{solution}

\begin{solution}{exo:g12:ka-and-pka:10}
$[\ce{H3O+}] = 10^{-5.6} = \qty{2.5e-6}{mol/L}$, शुद्ध पानी
(\qty{1.0e-7}{mol/L}) से 25 गुना अधिक।
\end{solution}

\begin{solution}{exo:g12:ka-and-pka:11}
$K_a K_b = \dfrac{[\ce{NH3}][\ce{H3O+}]}{[\ce{NH4+}]c^\circ} \times
\dfrac{[\ce{NH4+}][\ce{OH-}]}{[\ce{NH3}]c^\circ} = K_e$, इसलिए
$pK_a = pK_e - pK_b = 14.0 - 4.74 = 9.26$।
\end{solution}

\begin{solution}{exo:g12:ka-and-pka:12}
\qty{0.0010}{mol/L} पर: $x^2 + \num{1.74e-5}\,x - \num{1.74e-8} = 0$,
$x = \qty{1.24e-4}{mol/L}$, \omterm{def:g9:acids-bases-ph:ph}{pH} 3.91: 0.52 की वृद्धि, 1 की नहीं। तनु होने पर
\omterm{def:g12:ka-and-pka:strong-acid}{दुर्बल अम्ल} पानी से अधिक अनुपात में अभिक्रिया करता है (यहाँ
\qty{4}{\percent} के मुक़ाबले \qty{12}{\percent}), जो \omterm{def:g10:concentration-and-dilution:dilution}{तनुकरण} की
आंशिक भरपाई कर देता है।
\end{solution}

\begin{solution}{exo:g12:ka-and-pka:13}
अम्ल के रूप में: \ce{HCO3- + H2O <=> CO3^{2-} + H3O+}, $K = 10^{-10.33}$।
क्षारक के रूप में: \ce{HCO3- + H2O <=> H2CO3 + OH-} (\ce{H2CO3} यहाँ \ce{CO2,H2O} को दर्शाता है), $K = K_e/K_a = 10^{-14.0+6.37}
= 10^{-7.63}$। क्षारक अभिक्रिया का स्थिरांक बड़ा है: विलयन
थोड़ा क्षारकीय है।
\end{solution}

\begin{solution}{exo:g12:ka-and-pka:14}
कोई अम्ल, कितना भी तनु हो, विलयन को क्षारकीय नहीं बना सकता। जब अम्ल के
अपने \omterm{def:g12:ka-and-pka:oxonium}{ऑक्सोनियम आयन} पानी द्वारा दिए जाने वाले \qty{1e-7}{mol/L} से कम हो जाते हैं,
तो पानी के \omterm{def:g9:ions:ion}{आयन} हावी हो जाते हैं: \omterm{def:g9:acids-bases-ph:ph}{pH} नीचे से 7 की ओर बढ़ता है और
उसे कभी पार नहीं करता।
\end{solution}

\begin{solution}{exo:g12:ka-and-pka:15}
$n = 0.500 / 176.0 = \qty{2.84e-3}{mol}$, $c = \qty{0.0142}{mol/L}$;
$K_a = \num{9.12e-5}$; $x^2 + \num{9.12e-5}\,x - \num{1.30e-6} = 0$ से
$x = \qty{1.09e-3}{mol/L}$ मिलता है: \omterm{def:g9:acids-bases-ph:ph}{pH} 2.96।
\end{solution}

\begin{solution}{pb:g12:ka-and-pka:1}
\textbf{1.} \ce{H-C(=O)-O-H}: दिया जाने वाला हाइड्रोजन \ce{O-H} का है।

\textbf{2.} \ce{HCOOH}/\ce{HCOO-}; \ce{HCOOH + H2O <=> HCOO- + H3O+}।

\textbf{3.} $K_a = \dfrac{[\ce{HCOO-}][\ce{H3O+}]}{[\ce{HCOOH}]c^\circ} =
10^{-3.74} = \num{1.82e-4}$।

\textbf{4.} एथेनोइक अम्ल से नीचे ($3.74 < 4.76$): अधिक प्रबल।

\textbf{5.} $M = \qty{46.0}{g/mol}$: $0.010 \times 46.0 = \qty{0.46}{g}$।

\textbf{6.} \ce{HCOOH}: $0.010 - x$; \ce{HCOO-}: $x$; \ce{H3O+}: $x$।

\textbf{7.} $\dfrac{x^2}{0.010 - x} = \num{1.82e-4}$।

\textbf{8.} $x^2 + \num{1.82e-4}\,x - \num{1.82e-6} = 0$:
$x = \qty{1.26e-3}{mol/L}$।

\textbf{9.} $\mathrm{pH} = -\log(\num{1.26e-3}) = 2.90$।

\textbf{10.} $\num{1.26e-3}/0.010 = \qty{12.6}{\percent}$।

\textbf{11.} \qty{1.26e-3}{mol/L}, \qty{1e-7}{mol/L} के दस हज़ार गुने से
भी अधिक है: पानी के अपने \omterm{def:g9:ions:ion}{आयन} नगण्य हैं।

\textbf{12.} \omterm{def:g9:acids-bases-ph:ph}{pH} 2.0।

\textbf{13.} $0.010 / \num{1.26e-3} \approx 8$ गुना अधिक।

\textbf{14.} हाइड्रोक्लोरिक अम्ल एक \omterm{def:g12:ka-and-pka:strong-acid}{प्रबल अम्ल} है, जो पूरी तरह \omterm{def:g9:ions:ion}{आयनों} में
बदल जाता है; मेथेनोइक अम्ल दुर्बल है: पानी के साथ उसकी अभिक्रिया एक
साम्य पर रुक जाती है जहाँ उसका अधिकांश \ce{HCOOH} के रूप में बचा रहता है।

\textbf{15.} \ce{HCOOH + HCO3- -> HCOO- + CO2 + H2O}।

\textbf{16.} इस व्यंजक
\[
  K = \frac{[\ce{HCOO-}]\,[\ce{CO2}]}{[\ce{HCOOH}]\,[\ce{HCO3-}]}
\]
के अंश और हर को $[\ce{H3O+}]$ से गुणा करने पर $K = K_a(\ce{HCOOH}) / K_a(\ce{CO2,H2O}) =
\num{1.82e-4}/\num{4.3e-7} \approx 420$ मिलता है।

\textbf{17.} हाँ, $K \gg 1$; कार्बन डाइऑक्साइड बुलबुलों के रूप में निकलती है।

\textbf{18.} हाइड्रॉक्साइड एक \omterm{def:g12:ka-and-pka:strong-acid}{प्रबल क्षारक} है और स्वयं त्वचा पर आक्रमण करता;
हाइड्रोजनकार्बोनेट एक \omterm{def:g12:ka-and-pka:strong-acid}{दुर्बल क्षारक} है जो अम्ल को उदासीन करता है और
किसी भी अधिकता में सौम्य रहता है।

\textbf{19.} \omterm{def:g9:acids-bases-ph:ph}{pH} 2.9।
\end{solution}
